\section{HJB 方程与动态规划}
\label{sec:16-Control-and-Reinforcement-Learning/01-Optimal-Control/03}

\subsection{轨线有了，一阵风过后呢}

你控制一架无人机。状态 \(x\in\R^{12}\) 记着位置 3 维、速度 3 维、姿态角 3 维、角速度 3 维。控制 \(u\in\R^4\) 是四个旋翼的推力。动力学 \(\dot x = f(x,u)\) 有完整的空气动力学模型，代价函数也写清楚了：偏离目标轨迹罚 \(Q\)，控制用力过猛罚 \(R\)，终点偏差罚 \(S\)。

先按\hyperref[sec:16-Control-and-Reinforcement-Learning/01-Optimal-Control/02]{最大值原理}办。给定初值 \(x(0)\)，写 Hamiltonian \(H = L + \lambda^{\trans}f\)，得到耦合方程组 \(\dot x = \partial H/\partial\lambda\)、\(\dot\lambda = -\partial H/\partial x\)，外加逐点最小化 \(u^* = \argmin_u H\)。这是一个两点边值问题：\(x\) 有初值，\(\lambda\) 的终端条件从终值代价 \(\Phi\) 来。打靶法跑几轮，最优轨线和配套的控制序列都拿到了。

飞到 \(t=3.7\) 秒，一阵侧风推过来。位置偏了 0.4 米，速度方向也变了。你手上那条控制序列是用\emph{旧初值}算的——新的当前状态并不在那条轨线上。把原序列从新状态接着跑？不是最优。要拿到新状态出发的最优控制，你得把十二维 BVP 重解一遍。

每 10 毫秒跑一次的控制循环里重解 BVP？算不完。你缺的不是更快的求解器。你需要另一种形态的解：一个函数 \(u^*(t,x)\)，输入当前时刻和当前状态，立刻吐出最优控制。不管系统被风推到哪个状态，直接查答案，不用回头跑优化。

\subsection{值函数：把未来压进一个数}

先定义清楚目标。从时刻 \(t\)、状态 \(x\) 出发，之后全部采取最优控制，最小剩余代价是多少：

\[
V(t,x) = \inf_{u(\cdot)}
\Bigl\{
\Phi(x(T)) + \int_t^T L(x(s),u(s))\,ds
\Bigr\},
\]

其中 \(\dot x = f(x,u)\) 约束着整条积分路径，终点 \(x(T)\) 是你选的 \(u(\cdot)\) 从 \((t,x)\) 实际推到的地方。终点边界条件：

\[
V(T,x) = \Phi(x),
\]

到了终点不再有运行过程，只剩终值代价。

这个定义里还没有算法。它只是一个下确界——如果你能遍历全部控制函数，这是你能做到的最好结果。任何具体策略给出的代价顶多是 \(V(t,x)\) 的一个上界，永远打不破它。

手机导航的``预计剩余时间''是个接近的日常版本。它根据当前所在路口、此刻路况，估计从\emph{这里}到目的地的最短耗时。你已经开了多久少有单独立项的意义——它们已沉淀为当前位置了。原定路线突然封路，旧计划过期，但当前状态（位置、时间、路况）仍包含继续决策所需的所有信息，从新状态重新查值函数就好。

想深一层。两个位置眼下都等一个红灯，各预算一分钟。一个后面接畅通主路，另一个后面是持续拥堵。当前一分钟代价相同，值函数却不同。值函数比的是\emph{完整的剩余未来}，不是眼前这一步谁更快。这也是它和沿具体轨线定义的代价 \(J\) 的根本差别：\(J\) 针对一条具体的控制函数，哪怕那条控制不是最优的，\(J\) 总能算出一个数；\(V\) 是对全部可能的控制取最值，测量的是\emph{那个位置本身的潜力}。

\subsection{Bellman 最优性原理：尾巴不能有更优的替法}

思路是这样的。从 \((t,x)\) 出发的最优策略，先跑极短的一小段 \(h\)，到达 \(x(t+h)\)。从那个新状态往后，剩余控制序列不可能有比最优更优的——否则把那段更优的尾巴换上去，整条策略的代价就比原来的``最优''还低，自相矛盾：

\[
V(t,x) =
\inf_{u\text{ 在 }[t,t+h]}
\Bigl\{
\int_t^{t+h} L(x(s),u(s))\,ds
+ V\bigl(t+h,x(t+h)\bigr)
\Bigr\}.
\]

这就是 Bellman 最优性原理：全局最优路径的每一截尾部，从各自起点的视角看也必须最优。

它提供的是一个自洽条件，不是一个计算公式。和上一节的最大原理比着看差异：最大原理通过扰动整条轨线做一阶变分，得到一组必要条件；Bellman 原理通过切掉最开头一小截，得到一个递归结构。两条路从不同方向逼近同一个最优解，沿最优轨线给出等价的信息。

\subsection{\(h\to0\)：从递归到偏微分方程}

假设 \(V\) 光滑。在 \([t,t+h]\) 上把控制固定为常值 \(u\)，一阶展开：

\[
x(t+h) = x + h f(x,u) + o(h),
\qquad
\int_t^{t+h} L\,ds = h L(x,u) + o(h).
\]

\[
V(t+h,x(t+h)) = V(t,x) + h\partial_t V + h\nabla_x V\cdot f(x,u) + o(h).
\]

代回 Bellman 等式，消去两边的 \(V(t,x)\)，除以 \(h\)，令 \(h\to0\)：

\[
-\partial_t V(t,x) = \min_{u\in U}
\Bigl\{
L(x,u) + \nabla_x V(t,x)\cdot f(x,u)
\Bigr\},
\qquad
V(T,x) = \Phi(x).
\]

这就是 Hamilton--Jacobi--Bellman 方程。一阶非线性 PDE，空间 \(n+1\) 维（\(n\) 个状态加时间），时间方向从 \(T\) 倒着往 \(0\) 走。一旦你解出了 \(V\)，最优反馈律直接从内部的逐点最小化读出来：

\[
u^*(t,x) \in \argmin_{u\in U}
\Bigl\{
L(x,u) + \nabla_x V(t,x)\cdot f(x,u)
\Bigr\}.
\]

注意 \(\nabla_x V\) 的含义：它衡量``当前状态变好一单位，未来代价能省多少''。这正是上一节协态 \(\lambda\) 扮演的角色。沿最优轨线取 \(\lambda(t) = \nabla_x V(t,x^*(t))\)，HJB 方程对 \(x\) 求导就能恢复伴随方程 \(\dot\lambda = -\nabla_x H\)。

两条理论路线在这里交汇了——一条从全局值函数出发，一条从单条轨线的一阶必要条件出发。交汇处是同一个最优解，但代价的形态不同：最大值原理解一个 \(2n\) 维 BVP，HJB 解一个 \(n+1\) 维非线性 PDE 外加逐点最小化。不存在免费的午餐，交换的是计算负担的类型。开环解单点但扛不住扰动；反馈解全局但维数压人。

\subsection{LQR：二次假设把 PDE 压成 ODE}

退到恰好能完整算到底的情形。线性二次调节器（LQR）：

\[
\dot x = Ax + Bu,
\]

\[
J = \frac12 x(T)^{\trans}S x(T) + \frac12\int_0^T\bigl(x^{\trans}Q x + u^{\trans}R u\bigr)dt,
\]

\(Q,S\succeq 0\)，\(R\succ 0\)。偏离和用力都按二次收费。

猜值函数也是二次型：\(V(t,x) = \frac12 x^{\trans}P(t)x\)，\(P(t)\) 对称，待定。于是 \(\nabla_x V = P(t)x\)。HJB 右端关于 \(u\) 的部分：

\[
\frac12 u^{\trans}R u + x^{\trans}P B u.
\]

对 \(u\) 求梯度并置零（\(\frac12 u^{\trans}R u\) 的梯度是 \(R u\)，\(x^{\trans}P B u\) 对 \(u\) 的梯度是 \(B^{\trans}P x\)）：

\[
R u + B^{\trans}P x = 0,\qquad
u^*(t,x) = -R^{-1}B^{\trans}P(t)\,x.
\]

最优控制自动成为状态的线性反馈。增益 \(K(t)=R^{-1}B^{\trans}P(t)\)。不需要每次做逐点最小化——代数结构已被锁定，问题只剩下确定 \(P(t)\)。

代回 HJB。左端：\(-\frac12 x^{\trans}\dot P x\)。右端展开：

\[
\frac12 x^{\trans}Q x + \frac12 x^{\trans}P B R^{-1}B^{\trans}P x
+ \frac12 x^{\trans}(A^{\trans}P+PA)x
- x^{\trans}P B R^{-1}B^{\trans}P x.
\]

最后两项合并：\(x^{\trans}P A x = \frac12 x^{\trans}(A^{\trans}P+PA)x\)（取对称部分），而 \(x^{\trans}P B R^{-1}B^{\trans}P x\) 的系数为 \(\frac12-1=-\frac12\)。消去二次型：

\[
-\dot P = Q + A^{\trans}P + PA - P B R^{-1}B^{\trans}P,\qquad
P(T)=S.
\]

Riccati 微分方程。原来跨越 \((t,x)\) 的非线性 PDE，被二次假设压成了只依赖 \(t\) 的矩阵 ODE。\(P\) 是 \(n\times n\) 对称矩阵，含 \(n(n+1)/2\) 个独立标量方程。\(n=12\) 的无人机是 78 个 ODE，从 \(T\) 往 0 倒着积分，几毫秒跑完。

运算分两步。离线：从 \(P(T)=S\) 倒积 Riccati 方程到 \(t=0\)，存增益 \(K(t)=R^{-1}B^{\trans}P(t)\)。在线：任意时刻、任意状态，\(u^* = -K(t)x\) 就是一个矩阵乘向量。被风吹偏了？状态 \(x\) 变了，用新 \(x\) 重新乘一次 \(K(t)\) 即可。反馈律以你采样状态的频率提供新控制，永远不需要回头跑优化。

\subsection{标量例子：权衡被压成两个数}

取最简系统 \(\dot x = u\)，无穷时域：

\[
J = \frac12\int_0^\infty \bigl(q x(t)^2 + r u(t)^2\bigr)dt,\qquad q>0,\ r>0.
\]

稳态下值函数与 \(t\) 无关：\(V(x)=\frac12 P x^2\)，\(P\ge0\) 是标量。稳态 HJB：

\[
0 = \min_u\Bigl\{\frac12 q x^2 + \frac12 r u^2 + P x u\Bigr\}.
\]

最小化：\(r u + P x = 0\)，即 \(u^* = -(P/r)x\)。代回：

\[
\frac12 q x^2 - \frac{P^2}{2r}x^2 = 0\;\Longrightarrow\; q-\frac{P^2}{r}=0.
\]

取非负根 \(P=\sqrt{qr}\)。于是

\[
u^* = -\sqrt{\frac{q}{r}}\,x,\qquad
V(x)=\frac12\sqrt{qr}\,x^2.
\]

两个数把整个权衡关系铸了进去。增大 \(q\)（偏离更不能忍）：\(\sqrt{q/r}\) 变大，控制更激进，\(V(x)\) 也变大——犯错变贵了。增大 \(r\)（更心疼费电）：\(\sqrt{q/r}\) 变小，控制更克制，但 \(V(x)\) 同样变大——省控制的代价是忍更大的偏离。

LQR 比手工调 PID 多了这一层建模依据。你先把``偏离一单位''和``多用一单位控制''的相对权重说清楚，Riccati 方程随后推出最优增益。权重本身仍然是你的建模判断——数学不替你判哪个更贵，但它保证：一旦做了判断，给出的反馈律在所有策略中使二次代价达到绝对最小。

\subsection{二维例子：位置和速度的联合反馈}

让系统稍微丰富一层。双积分器 \(\ddot x = u\)，写成状态空间：

\[
\dot x_1 = x_2,\qquad \dot x_2 = u,\qquad
J = \frac12\int_0^\infty\bigl(q x_1^2 + r u^2\bigr)dt.
\]

希望位置 \(x_1\) 趋近于零，代价中不直接罚速度（但位置偏离大时速度会跟着偏离，间接被罚）。矩阵形式：

\[
A=\begin{bmatrix}0&1\\0&0\end{bmatrix},\;
B=\begin{bmatrix}0\\1\end{bmatrix},\;
Q=\begin{bmatrix}q&0\\0&0\end{bmatrix},\;
R=[r].
\]

稳态 Riccati：\(0 = Q + A^{\trans}P + PA - PBR^{-1}B^{\trans}P\)。设 \(P=\begin{bmatrix}p_{11}&p_{12}\\p_{12}&p_{22}\end{bmatrix}\)，逐元素展开：

\[
\begin{bmatrix}q&0\\0&0\end{bmatrix}
+ \begin{bmatrix}0&0\\p_{11}&p_{12}\end{bmatrix}
+ \begin{bmatrix}0&p_{11}\\0&p_{12}\end{bmatrix}
- \frac1r\begin{bmatrix}p_{12}^2&p_{12}p_{22}\\p_{12}p_{22}&p_{22}^2\end{bmatrix}
= \begin{bmatrix}0&0\\0&0\end{bmatrix}.
\]

逐个方程：\(q-p_{12}^2/r=0 \Rightarrow p_{12}=\sqrt{qr}\)。\(p_{11}-p_{12}p_{22}/r=0\)（上三角项）。\(2p_{12}-p_{22}^2/r=0 \Rightarrow p_{22}=\sqrt{2p_{12}r}\)。

最终增益：

\[
K = R^{-1}B^{\trans}P = \frac1r\begin{bmatrix}p_{12}&p_{22}\end{bmatrix}
= \begin{bmatrix}\sqrt{\frac{q}{r}}&\sqrt{2}\bigl(\frac{q}{r}\bigr)^{1/4}\end{bmatrix}.
\]

最优控制 \(u^* = -Kx = -\sqrt{q/r}\,x_1 - \sqrt{2}(q/r)^{1/4}\,x_2\)。两项各有职责：位置偏离触发纠正，速度过快触发阻尼。LQR 一次输出了两者的增益\emph{和}比例——你不需要分别调两个旋钮再猜耦合效应，Riccati 方程直接给出完整的协调方案。

\subsection{光滑性的裂缝：HJB 不总可微}

整个推导假定了 \(V\) 光滑可微。但最优控制问题的解经常在策略切换处出现折角——最短时间控制的 bang-bang 切换就是典型：最优控制从最大正向跳到最大反向，值函数在切换面上拐了一个尖。折角处偏导数不存在，经典 HJB 不能逐点套用。

黏性解框架处理的就是这个问题。在不可微点，用局部光滑测试函数从上或从下方接触 \(V\)，把等式约束松弛为不等式约束。黏性解在 \(V\) 光滑时退化为经典解，在不可微时保留唯一性，且能被单调有限差分格式稳定逼近。

LQR 碰不上这个问题——二次型天生光滑。但一旦你的问题加了控制饱和（\(u_{\min}\le u\le u_{\max}\)）或状态约束（避开障碍），逐点最小化就变成了约束优化，值函数不再是单个二次型，HJB 也不能再降阶为 Riccati。黏性解就是理解这一类``解''的内涵时必须进入的语言。

\subsection{维数灾难：画一张全状态地图的代价}

HJB 方案要在状态空间每个点存一个值。\(d\) 维空间，每维取 \(m\) 个格点，一共 \(m^d\) 个格点。每个格点要存 \(V\)、跑最小化、与邻居耦合。取 \(m=100\)、\(d=6\)：\(10^{12}\) 个格点，远超任何计算资源。

这不是方法的问题——动态规划精确承接了``同时为全部可能状态求最优代价''的计算成本。问题只在于这张地图在大多数高维问题中画不起。

LQR 的价值不在泛用，而在精确识别了一个计算上可行的子空间：线性动力学 + 二次代价 + 无约束控制。在这子空间里，值函数被 \(n\times n\) 对称矩阵完整参数化，存储从 \(O(m^d)\) 坍缩到 \(O(n^2)\)。

三条路摆在一起：最大值原理绕过状态网格求单条轨线（不画全地图，丢了反馈）；HJB 画全地图（拿到反馈，承担维数爆炸）；LQR 同时拿到反馈和可计算性——代价是三个强假设：线性、二次、无约束。

卸下其中一个假设，工具就得跟着换。下一个单元拿掉``模型已知''，进入强化学习：不能对转移函数做精确期望，就从一条条采样中逼近值函数。Bellman 递归的结构还在，只是计算的载体从精确积分变成了随机采样和函数逼近的组合。
